% ============================================================
%  GUÍA 1: RAZONAMIENTO LÓGICO Y LENGUAJE MATEMÁTICO
%  Curso Introductorio — 3er Año
%  Autor: Gabriel Astudillo
%  Sesiones cubiertas: 1, 2, 3
% ============================================================

\documentclass[12pt, a4paper]{article}

% ── Paquetes ─────────────────────────────────────────────────

\usepackage{mathwizards-guia}
\mwguia{Guía 1 — Razonamiento Lógico y Lenguaje Matemático}{Gabriel Astudillo $\cdot$ Curso 3er Año}

\begin{document}
% ============================================================

% ── Portada / Título ─────────────────────────────────────────
\mwportada{Guía de Estudio 1}{Razonamiento Lógico y Lenguaje Matemático}{Lógica Proposicional, Método de Pólya y Traducción Algebraica}

\vspace{4pt}
\begin{cuadroinfo}
\small
\textbf{Presentación:} Bienvenido al inicio de un viaje por la ciencia. Esta primera guía está diseñada para que aprendas a \emph{pensar con orden}. Antes de lanzarnos a resolver ecuaciones o medir magnitudes, necesitamos afinar nuestra herramienta más importante: la mente. Aquí descubrirás cómo funciona el razonamiento lógico, cómo los matemáticos atacan problemas paso a paso, y cómo traducir situaciones cotidianas al lenguaje preciso del álgebra.
\\[4pt]
\textbf{Nivel de Dificultad:}\quad
\facil{} Básico / Directo \quad
\medio{} Intermedio \quad
\dificil{} Desafiante
\end{cuadroinfo}

\vspace{8pt}

% ============================================================
\section{Introducción a la Lógica Proposicional}
% ============================================================

\subsection{\textquestiondown Qué es una proposición?}

Imagina que estás en un juicio y alguien debe declarar si algo es \emph{verdadero} o \emph{falso}. No puede responder «tal vez» ni «depende»: la respuesta es una u otra. Ese tipo de enunciado, que obligatoriamente es verdadero o falso (nunca ambos, nunca ninguno), se llama \textbf{proposición}.

\begin{cuadroteoria}
\textbf{Definición:} Una \textbf{proposición} es un enunciado declarativo al que se le puede asignar exactamente uno de dos \textbf{valores de verdad}: Verdadero (\textbf{V}) o Falso (\textbf{F}).
\end{cuadroteoria}

\noindent\textbf{Ejemplos de proposiciones:}
\begin{itemize}[leftmargin=*]
  \item «La Tierra gira alrededor del Sol.» → \textbf{V}
  \item «$2 + 3 = 7$» → \textbf{F}
  \item «Caracas es la capital de Venezuela.» → \textbf{V}
  \item «El agua hierve a 50\,°C a nivel del mar.» → \textbf{F}
\end{itemize}

\noindent\textbf{Ejemplos de enunciados que NO son proposiciones:}
\begin{itemize}[leftmargin=*]
  \item «¿Qué hora es?» → Es una pregunta, no afirma nada.
  \item «¡Cierra la puerta!» → Es una orden.
  \item «$x + 5 = 12$» → Depende del valor de $x$; no podemos asignarle V o F sin más información.
  \item «Esta oración es falsa.» → Genera una paradoja (si es V entonces es F, y viceversa).
\end{itemize}

\subsection{Proposiciones atómicas y moleculares}

Las proposiciones pueden ser simples o compuestas:

\begin{itemize}[leftmargin=*]
  \item Una \textbf{proposición atómica} (o simple) no se puede descomponer en proposiciones más pequeñas. Ejemplo: «Está lloviendo.»
  \item Una \textbf{proposición molecular} (o compuesta) se forma al unir dos o más proposiciones atómicas mediante palabras especiales llamadas \textbf{conectivos lógicos}. Ejemplo: «Está lloviendo \emph{y} hace frío.»
\end{itemize}

Para trabajar cómodamente, representamos cada proposición atómica con una letra minúscula: $p$, $q$, $r$, etc.

\subsection{Los conectivos lógicos}

Los conectivos lógicos son las «palabras pegamento» que unen proposiciones. Veamos los tres fundamentales:

\begin{cuadrosolucion}
\begin{enumerate}[label=\arabic*., leftmargin=*]
  \item \textbf{Conjunción} (símbolo: $\wedge$, se lee «\emph{y}»): La proposición $p \wedge q$ es verdadera \textbf{solo cuando ambas} $p$ y $q$ son verdaderas.
  
  \textit{Ejemplo:} «Aprobé matemática \textbf{y} aprobé física.» Solo es verdad si aprobaste las dos.

  \item \textbf{Disyunción} (símbolo: $\vee$, se lee «\emph{o}»): La proposición $p \vee q$ es verdadera \textbf{cuando al menos una} de las dos es verdadera.
  
  \textit{Ejemplo:} «Voy al cine \textbf{o} me quedo en casa.» Es verdad si haces una cosa, la otra, o incluso ambas.

  \item \textbf{Condicional} (símbolo: $\rightarrow$, se lee «\emph{si \ldots\ entonces}»): La proposición $p \rightarrow q$ es falsa \textbf{solamente} cuando $p$ es verdadera y $q$ es falsa. En todos los demás casos es verdadera.
  
  \textit{Ejemplo:} «\textbf{Si} estudias, \textbf{entonces} aprobarás.» Solo puedes decir que esta promesa se rompió si efectivamente estudiaste y aun así no aprobaste.
\end{enumerate}
\end{cuadrosolucion}

\subsection{Tablas de verdad}

Una \textbf{tabla de verdad} es simplemente un cuadro donde listamos todas las combinaciones posibles de valores de verdad de las proposiciones atómicas y calculamos el resultado de la proposición compuesta. Veamos las tablas para los tres conectivos:

\begin{center}
\begin{tabular}{cc|c|c|c}
\toprule
$p$ & $q$ & $p \wedge q$ & $p \vee q$ & $p \rightarrow q$ \\
\midrule
V & V & V & V & V \\
V & F & F & V & F \\
F & V & F & V & V \\
F & F & F & F & V \\
\bottomrule
\end{tabular}
\end{center}

\begin{cuadroadvertencia}
\textbf{Dato curioso:} La línea que más sorprende a los estudiantes es la última del condicional: cuando $p$ es falsa, $p \rightarrow q$ siempre es verdadera, sin importar $q$. Piénsalo así: si tu profesor te dice «si sacas 20 puntos, te regalo un libro» y tú no sacas 20, el profesor no rompió su promesa, pase lo que pase. Una promesa solo se rompe si se cumple la condición y no se cumple lo prometido.
\end{cuadroadvertencia}

\subsection{Traducción del lenguaje natural a la lógica simbólica}

Traducir del español a símbolos lógicos consiste en:
\begin{enumerate}
  \item \textbf{Identificar} las proposiciones atómicas y asignarles letras.
  \item \textbf{Detectar} los conectivos lógicos en el enunciado.
  \item \textbf{Escribir} la expresión simbólica.
\end{enumerate}

\noindent\textbf{Ejemplo:} «Si llueve y no tengo paraguas, entonces me mojaré.»
\begin{itemize}
  \item $p$: «Llueve.»
  \item $q$: «Tengo paraguas.»
  \item $r$: «Me mojaré.»
  \item Traducción: $(p \wedge \neg q) \rightarrow r$
\end{itemize}

Observa que la palabra «no» se representa con el símbolo $\neg$ (\textbf{negación}), que simplemente invierte el valor de verdad de una proposición.

\newpage

% ============================================================
\subsection{Ejercicios: Lógica Proposicional}
% ============================================================

\begin{enumerate}[leftmargin=*]
  \item \facil{} Un detective llega a la escena del crimen y encuentra tres sospechosos: Ana, Bruno y Carlos. Se sabe que:
  \begin{itemize}
    \item Si Ana es culpable, entonces Bruno también lo es.
    \item Bruno es inocente.
  \end{itemize}
  ¿Puedes concluir algo sobre Ana? Justifica tu razonamiento.

  \item \facil{} En una isla viven dos tipos de habitantes: los \emph{caballeros} (siempre dicen la verdad) y los \emph{pícaros} (siempre mienten). Te encuentras a un habitante que dice: «Yo soy un pícaro.» ¿Es caballero o pícaro? ¿Por qué?

  \item \facil{} Determina cuáles de los siguientes enunciados son proposiciones y asígnales un valor de verdad (V o F) cuando sea posible:
  \begin{enumerate}[label=\alph*)]
    \item «El número 15 es par.»
    \item «¡Pásame el borrador!»
    \item «$3 \times 4 = 12$»
    \item «¿Venezuela queda en Europa?»
    \item «El oxígeno es un gas a temperatura ambiente.»
  \end{enumerate}

  \item \facil{} Traduce al lenguaje simbólico, definiendo las proposiciones atómicas que necesites:
  \begin{enumerate}[label=\alph*)]
    \item «Hace calor y estoy cansado.»
    \item «Si estudio, entonces apruebo el examen.»
    \item «Llueve o hace sol.»
  \end{enumerate}

  \item \medio{} Un profesor afirma: «Si entregas la tarea a tiempo \textbf{y} participas en clase, entonces tendrás puntos extra.» Marcos entregó la tarea a tiempo pero no participó en clase, y no obtuvo puntos extra. ¿El profesor cumplió su promesa o la rompió? Justifica usando el condicional.

  \item \medio{} Construye la tabla de verdad completa para la proposición $(p \vee q) \rightarrow p$.

  \item \medio{} En un juego de cartas, se establece la siguiente regla: «Si una carta tiene una vocal en un lado, entonces tiene un número par en el otro.» Tienes cuatro cartas sobre la mesa mostrando: \textbf{A}, \textbf{K}, \textbf{4}, \textbf{7}. ¿Cuáles cartas debes voltear para verificar si la regla se cumple? Justifica tu respuesta lógicamente.

  \item \medio{} Tres amigos hacen las siguientes afirmaciones:
  \begin{itemize}
    \item Luis: «Si María fue al cine, entonces Pedro se quedó en casa.»
    \item María: «Yo fui al cine.»
    \item Pedro: «Yo no me quedé en casa.»
  \end{itemize}
  Si todos dicen la verdad excepto uno, ¿quién miente?

  \item \medio{} Traduce al lenguaje simbólico y luego construye la tabla de verdad: «Si no llueve, entonces iré al parque o me quedaré leyendo.»

  \item \dificil{} En la isla de caballeros y pícaros, te encuentras con dos habitantes, $A$ y $B$. El habitante $A$ dice: «Al menos uno de nosotros es pícaro.» Determina qué tipo es cada uno.

  \item \dificil{} Un acertijo clásico: Tienes tres cajas etiquetadas «Manzanas», «Naranjas» y «Mixta», pero \textbf{todas las etiquetas están mal colocadas}. Solo puedes sacar \textbf{una fruta} de \textbf{una caja} para determinar el contenido correcto de las tres cajas. ¿De qué caja sacas la fruta y por qué? Explica tu razonamiento lógico paso a paso.

  \item \dificil{} Considera la afirmación: «Si un número es divisible por 6, entonces es divisible por 2 \textbf{y} por 3.» 
  \begin{enumerate}[label=\alph*)]
    \item Escríbela en lenguaje simbólico.
    \item ¿Es verdadera o falsa? Justifica.
    \item Escribe la afirmación \emph{recíproca} (intercambiando la hipótesis con la conclusión). ¿Es verdadera? Da un ejemplo o contraejemplo.
  \end{enumerate}
\end{enumerate}

\newpage

% ============================================================
\section{El Método de George Pólya}
% ============================================================

\subsection{\textquestiondown Cómo resuelven los problemas los buenos matemáticos?}

Seguramente te ha pasado: lees un problema, no sabes por dónde empezar, y terminas dejándolo en blanco. La buena noticia es que resolver problemas \textbf{es una habilidad que se puede aprender}. George Pólya, un famoso matemático húngaro, dedicó años a observar cómo las personas resuelven problemas y resumió su método en cuatro fases claras.

\begin{cuadroteoria}
\textbf{El Método de Pólya en 4 fases:}
\begin{enumerate}[label=\textbf{Fase \arabic*:}, leftmargin=*]
  \item \textbf{Comprender el problema} — ¿Qué te piden? ¿Qué datos tienes? ¿Qué condiciones hay? Reformula el problema con tus propias palabras.
  \item \textbf{Concebir un plan} — ¿Has visto algo parecido antes? ¿Puedes simplificarlo? ¿Hay un patrón? Busca analogías con problemas resueltos.
  \item \textbf{Ejecutar el plan} — Lleva a cabo tu estrategia paso a paso, con cuidado y orden. Si te atascas, vuelve a la Fase 2.
  \item \textbf{Visión retrospectiva} — ¿Tu resultado tiene sentido? ¿Puedes verificarlo? ¿Podrías haberlo resuelto de otra forma?
\end{enumerate}
\end{cuadroteoria}

\begin{figure}[h!]
  \centering
  \begin{tikzpicture}[
      phase/.style = {circle, draw=mwprimario, fill=mwprimario!12, thick, minimum size=3cm, text width=2.6cm, align=center, text=mwprimario, font=\bfseries\small},
      arrow/.style = {-{Latex[scale=1.2]}, thick, draw=mwprimario}
    ]
    \node (f1) [phase] at (90:2.8) {Fase 1:\\Comprender\\el problema};
    \node (f2) [phase] at (0:2.8) {Fase 2:\\Concebir\\un plan};
    \node (f3) [phase] at (270:2.8) {Fase 3:\\Ejecutar\\el plan};
    \node (f4) [phase] at (180:2.8) {Fase 4:\\Visión\\retrospectiva};

    \draw [arrow] (f1) to[bend left=20] (f2);
    \draw [arrow] (f2) to[bend left=20] (f3);
    \draw [arrow] (f3) to[bend left=20] (f4);
    \draw [arrow] (f4) to[bend left=20] (f1);
    
    \draw [arrow, dashed, mwrojonaranja, bend left=30] (f3) to node[midway, below left, text=mwrojonaranja, font=\scriptsize\bfseries] {Replanificar} (f2);
    \draw [arrow, dashed, mwrojonaranja, bend left=30] (f4) to node[midway, above left, text=mwrojonaranja, font=\scriptsize\bfseries] {Reevaluar} (f1);
  \end{tikzpicture}
  \caption{Las cuatro fases del Método de Pólya.}
  \label{fig:polya}
\end{figure}

\subsection{Ejemplo guiado con las 4 fases}

\noindent\textbf{Problema:} En un estacionamiento hay bicicletas y automóviles. En total se cuentan 18 vehículos y 50 ruedas. ¿Cuántas bicicletas y cuántos automóviles hay?

\begin{cuadrosolucion}
\textbf{Fase 1 — Comprender:} 
\begin{itemize}[leftmargin=*]
  \item \textit{¿Qué nos piden?} La cantidad de bicicletas y la cantidad de automóviles.
  \item \textit{¿Qué sabemos?} Hay 18 vehículos en total. Hay 50 ruedas en total. Las bicicletas tienen 2 ruedas y los autos tienen 4.
\end{itemize}
\end{cuadrosolucion}

\begin{cuadrosolucion}
\textbf{Fase 2 — Planificar:}
Podemos usar una estrategia de «suposición extrema»: ¿qué pasaría si todos fueran bicicletas? Tendríamos $18 \times 2 = 36$ ruedas, pero hay 50. Nos faltan $50 - 36 = 14$ ruedas. Cada auto que «reemplace» una bicicleta agrega 2 ruedas extra. Entonces, el número de autos sería $14 \div 2 = 7$.
\end{cuadrosolucion}

\begin{cuadrosolucion}
\textbf{Fase 3 — Ejecutar:}
\begin{itemize}
  \item Automóviles: $7$
  \item Bicicletas: $18 - 7 = 11$
\end{itemize}
\end{cuadrosolucion}

\begin{cuadrosolucion}
\textbf{Fase 4 — Verificar:}
$11$ bicicletas $\times$ 2 ruedas $= 22$ ruedas. $7$ autos $\times$ 4 ruedas $= 28$ ruedas. Total: $22 + 28 = 50$ ruedas. $\checkmark$ Total vehículos: $11 + 7 = 18$. $\checkmark$
\end{cuadrosolucion}

\newpage

% ============================================================
\subsection{Ejercicios: Método de Pólya}
% ============================================================
Para cada problema, aplica explícitamente las cuatro fases del Método de Pólya. Escribe cada fase por separado.

\begin{enumerate}[leftmargin=*, resume]
  \item \facil{} En una granja hay gallinas y conejos. Se cuentan 20 cabezas y 56 patas. ¿Cuántas gallinas y cuántos conejos hay?

  \item \facil{} María tiene el doble de edad que su hermano Juan. Dentro de 5 años, la suma de sus edades será 37 años. ¿Qué edad tiene cada uno ahora?

  \item \facil{} Un rectángulo tiene un perímetro de 40\,cm. Si el largo es 4\,cm mayor que el ancho, ¿cuáles son las dimensiones del rectángulo?

  \item \facil{} Tienes billetes de 5 y de 10 bolívares. En total son 12 billetes que suman 85 bolívares. ¿Cuántos billetes de cada tipo tienes?

  \item \medio{} Un tren sale de la ciudad A hacia la ciudad B a 80\,km/h. Dos horas después, otro tren sale de B hacia A a 100\,km/h. Si las ciudades están a 600\,km de distancia, ¿después de cuántas horas (desde que salió el segundo tren) se encuentran?

  \item \medio{} En una prueba de 25 preguntas, cada respuesta correcta vale 4 puntos y cada incorrecta resta 1 punto (no hay preguntas sin responder). Un estudiante obtuvo 60 puntos. ¿Cuántas respuestas correctas e incorrectas tuvo?

  \item \medio{} Una piscina se puede llenar con la llave A en 6 horas y con la llave B en 4 horas. Si se abren ambas llaves al mismo tiempo, ¿en cuántas horas se llena la piscina?

  \item \medio{} Tres amigos tienen juntos 120 caramelos. Pedro tiene el doble que Ana, y Ana tiene 10 más que Luis. ¿Cuántos caramelos tiene cada uno?

  \item \dificil{} Un caracol sube por una pared de 10 metros. De día sube 3 metros, pero de noche resbala 2 metros. ¿En cuántos días llega a la cima?

  \item \dificil{} En una competencia, los participantes reciben 5 puntos por cada prueba superada y pierden 3 puntos por cada prueba fallida. Después de 20 pruebas, un competidor tiene 52 puntos. ¿Cuántas pruebas superó?

  \item \dificil{} Tienes una balanza de dos platillos y tres monedas de apariencia idéntica, pero una de ellas es falsa y pesa menos que las otras. ¿Cómo puedes identificar la moneda falsa usando la balanza \textbf{una sola vez}?

  \item \dificil{} Un reloj da campanadas cada hora (1 campanada a la 1:00, 2 a las 2:00, ..., 12 a las 12:00). ¿Cuántas campanadas da en total desde la 1:00 hasta las 12:00 inclusive? Si además da 1 campanada cada media hora, ¿cuántas campanadas da en un día completo (24 horas)?
\end{enumerate}

\newpage

% ============================================================
\section{Del Lenguaje Natural al Lenguaje Algebraico}
% ============================================================

\subsection{Variables, constantes e incógnitas}

En la vida cotidiana usamos palabras para describir situaciones. En matemática, usamos \textbf{símbolos} para hacer lo mismo, pero con mayor precisión. Para ello, necesitamos entender tres conceptos clave:

\begin{cuadroteoria}
\begin{itemize}[leftmargin=*]
  \item \textbf{Constante:} Un valor fijo que no cambia. Ejemplo: el número 5, el precio de una entrada que cuesta 20\,Bs, la aceleración de la gravedad $g = 9{,}8\;\text{m/s}^2$.
  \item \textbf{Variable:} Un símbolo (generalmente una letra) que puede tomar diferentes valores. Ejemplo: la temperatura $T$ de un día cualquiera, la velocidad $v$ de un carro en movimiento.
  \item \textbf{Incógnita:} Una variable cuyo valor específico desconocemos y queremos encontrar. Ejemplo: «Un número sumado con 7 da 15. ¿Cuál es ese número?» → Aquí la incógnita es ese «número» misterioso, que representamos con $x$.
\end{itemize}
\end{cuadroteoria}

\subsection{Traducir enunciados a expresiones algebraicas}

La clave está en reconocer las \textbf{palabras clave} del español que corresponden a operaciones matemáticas:

\begin{center}
\small
\begin{tabular}{p{5.5cm} c p{4cm}}
\toprule
\textbf{Expresión en español} & \textbf{Operación} & \textbf{Símbolo} \\
\midrule
«la suma de», «más que», «aumentado en», «excede en» & Adición & $a + b$ \\
«la diferencia de», «menos que», «disminuido en» & Sustracción & $a - b$ \\
«el producto de», «multiplicado por», «el doble/triple de» & Multiplicación & $a \cdot b$, $2a$, $3a$ \\
«el cociente de», «dividido entre», «la mitad/tercera parte de» & División & $\frac{a}{b}$, $\frac{a}{2}$, $\frac{a}{3}$ \\
«el cuadrado de», «al cubo» & Potenciación & $a^2$, $a^3$ \\
«es igual a», «resulta», «da como resultado» & Igualdad & $=$ \\
\bottomrule
\end{tabular}
\end{center}

\subsection{Ejemplos de traducción}

\noindent\textbf{Ejemplo 1:} «El triple de un número, disminuido en 4, es igual a 17.»
\begin{itemize}
  \item «un número» → $x$ (incógnita)
  \item «el triple de un número» → $3x$
  \item «disminuido en 4» → $3x - 4$
  \item «es igual a 17» → $3x - 4 = 17$
\end{itemize}

\noindent\textbf{Ejemplo 2:} «La suma de dos números consecutivos es 45.»
\begin{itemize}
  \item Primer número: $n$
  \item El siguiente consecutivo: $n + 1$
  \item «la suma es 45» → $n + (n + 1) = 45$
\end{itemize}

\noindent\textbf{Ejemplo 3:} «El área de un rectángulo es el producto de su base por su altura.»
\begin{itemize}
  \item Área: $A$, base: $b$, altura: $h$
  \item Traducción: $A = b \cdot h$
\end{itemize}

\begin{cuadroadvertencia}
\textbf{Consejo:} Cuando traduzcas, léelo al revés para verificar. Si escribiste $3x - 4 = 17$, léelo como «el triple de $x$, menos 4, es igual a 17.» ¿Coincide con el enunciado original? ¡Perfecto!
\end{cuadroadvertencia}

\newpage

% ============================================================
\subsection{Ejercicios: Lenguaje Algebraico}
% ============================================================
Traduce cada enunciado a una expresión o ecuación algebraica. Define claramente tus variables.

\begin{enumerate}[leftmargin=*, resume]
  \item \facil{} «El doble de un número, aumentado en 9.»

  \item \facil{} «La mitad de un número es igual a 20.»

  \item \facil{} «La suma de tres números consecutivos.»

  \item \facil{} «Un número al cuadrado, menos el mismo número, da como resultado 12.»

  \item \medio{} «La edad de Pedro es el triple de la edad de su hijo. Juntos suman 48 años.» Escribe la ecuación que permite hallar la edad del hijo.

  \item \medio{} «El perímetro de un triángulo equilátero de lado $\ell$.»

  \item \medio{} «Un artículo cuesta $P$ bolívares. Si se le aplica un descuento del 15\,\%, ¿cuál es el precio final?»

  \item \medio{} «La velocidad de un móvil es igual a la distancia que recorre dividida entre el tiempo que tarda.» Escribe la fórmula correspondiente.

  \item \medio{} «Dos obreros pintan una cerca. El primero pinta 3 metros por hora y el segundo pinta 5 metros por hora. Juntos deben pintar una cerca de 40 metros.» Escribe la ecuación que modela el tiempo necesario.

  \item \dificil{} «El cuadrado de la suma de dos números es igual a la suma de sus cuadrados más el doble de su producto.» Escribe esta identidad algebraica.

  \item \dificil{} «Un tren recorre cierta distancia $d$. Si aumentara su velocidad en 20\,km/h, tardaría 1 hora menos en hacer el mismo recorrido.» Usando $v$ como la velocidad original y $d$ como la distancia, plantea dos ecuaciones que relacionen las variables.

  \item \dificil{} «En una tienda, un pantalón cuesta el doble que una camisa. Si compras 3 pantalones y 5 camisas, el total es 770\,Bs. Además, 2 camisas cuestan 40\,Bs más que un pantalón.» Plantea el sistema de ecuaciones completo (no lo resuelvas).
\end{enumerate}

\newpage

% ============================================================
\section{Clave de Respuestas}
% ============================================================

\subsection*{Lógica Proposicional (Ejercicios 1--12)}

\begin{enumerate}[leftmargin=*, label=\textbf{\arabic*.}]
  \item Ana es inocente. Si Ana fuera culpable, entonces Bruno sería culpable (por el condicional). Pero Bruno es inocente. Por contrapositiva ($\neg q \rightarrow \neg p$): si Bruno es inocente, entonces Ana también lo es.

  \item El habitante no puede ser pícaro (si lo fuera, estaría diciendo la verdad, lo cual contradice que los pícaros mienten). Tampoco puede ser caballero (si lo fuera, estaría afirmando algo falso). La situación es una \textbf{paradoja}: este enunciado no puede ser dicho por ningún habitante de la isla.

  \item a) Proposición, \textbf{F}. \quad b) No es proposición (orden). \quad c) Proposición, \textbf{V}. \quad d) No es proposición (pregunta). \quad e) Proposición, \textbf{V}.

  \item a) $p$: «Hace calor», $q$: «Estoy cansado» → $p \wedge q$. \quad b) $p$: «Estudio», $q$: «Apruebo el examen» → $p \rightarrow q$. \quad c) $p$: «Llueve», $q$: «Hace sol» → $p \vee q$.

  \item El profesor cumplió su promesa. La condición del condicional era «entregas la tarea a tiempo \textbf{y} participas en clase» ($p \wedge q$). Como Marcos no participó, $p \wedge q$ es falsa, y cuando la hipótesis es falsa, el condicional es verdadero independientemente de la conclusión.

  \item Tabla de verdad de $(p \vee q) \rightarrow p$:
  \begin{center}
  \small
  \begin{tabular}{cc|c|c}
  $p$ & $q$ & $p \vee q$ & $(p \vee q) \rightarrow p$ \\
  \hline
  V & V & V & V \\
  V & F & V & V \\
  F & V & V & F \\
  F & F & F & V \\
  \end{tabular}
  \end{center}

  \item Debes voltear \textbf{A} y \textbf{7}. La carta A tiene vocal, así que debes verificar que tenga número par detrás. La carta 7 es impar; si tiene vocal del otro lado, la regla se viola. La K y el 4 no requieren verificación (K no tiene vocal, y el 4 con cualquier letra cumple).

  \item Luis miente. Si María dice la verdad (fue al cine) y Pedro dice la verdad (no se quedó en casa), entonces la afirmación de Luis («Si María fue al cine, entonces Pedro se quedó en casa») es falsa, ya que la hipótesis es V y la conclusión es F.

  \item $p$: «Llueve», $q$: «Iré al parque», $r$: «Me quedaré leyendo». Traducción: $\neg p \rightarrow (q \vee r)$.
  \begin{center}
  \small
  \begin{tabular}{ccc|c|c|c}
  $p$ & $q$ & $r$ & $\neg p$ & $q \vee r$ & $\neg p \rightarrow (q \vee r)$ \\
  \hline
  V & V & V & F & V & V \\
  V & V & F & F & V & V \\
  V & F & V & F & V & V \\
  V & F & F & F & F & V \\
  F & V & V & V & V & V \\
  F & V & F & V & V & V \\
  F & F & V & V & V & V \\
  F & F & F & V & F & F \\
  \end{tabular}
  \end{center}

  \item $A$ es caballero y $B$ es pícaro. Si $A$ fuera pícaro, su afirmación «al menos uno de nosotros es pícaro» sería verdadera (porque $A$ mismo lo es), pero los pícaros mienten, contradicción. Entonces $A$ es caballero y su afirmación es verdadera: al menos uno es pícaro. Como $A$ no lo es, $B$ debe ser el pícaro.

  \item Sacas una fruta de la caja etiquetada «Mixta». Como la etiqueta es incorrecta, esa caja contiene solo un tipo de fruta. Si sacas una manzana, esa caja es «Manzanas». Entonces la caja etiquetada «Manzanas» (que está mal) no puede ser manzanas ni mixta (porque mixta ya se asignó... Espera, no). Corrección: la caja «Manzanas» tiene etiqueta incorrecta, y la caja «Naranjas» también. Como la caja real de manzanas ya fue identificada, la caja etiquetada «Naranjas» no puede ser naranjas (etiqueta incorrecta), así que es «Mixta», y la etiquetada «Manzanas» es «Naranjas».

  \item a) $p$: «un número es divisible por 6», $q$: «es divisible por 2», $r$: «es divisible por 3». Expresión: $p \rightarrow (q \wedge r)$. \quad b) Verdadera: todo múltiplo de 6 es múltiplo de 2 y de 3 por definición ($6 = 2 \times 3$). \quad c) Recíproca: $(q \wedge r) \rightarrow p$: «Si un número es divisible por 2 y por 3, entonces es divisible por 6.» También es verdadera, ya que $\text{mcm}(2,3) = 6$.
\end{enumerate}

\subsection*{Método de Pólya (Ejercicios 13--24)}

\begin{enumerate}[leftmargin=*, label=\textbf{\arabic*.}]
  \setcounter{enumi}{12}
  \item 12 gallinas y 8 conejos. (Si todas fueran gallinas: $20 \times 2 = 40$ patas. Faltan $56 - 40 = 16$. Cada conejo agrega 2 patas extra: $16 \div 2 = 8$ conejos, $20 - 8 = 12$ gallinas.)

  \item Juan tiene 9 años y María tiene 18 años. (Sea $x$ la edad de Juan. María: $2x$. Dentro de 5 años: $(x + 5) + (2x + 5) = 37 \Rightarrow 3x + 10 = 37 \Rightarrow x = 9$.)

  \item Ancho: 8\,cm, largo: 12\,cm. ($2(a + a + 4) = 40 \Rightarrow 2a + 4 = 20 \Rightarrow a = 8$.)

  \item 7 billetes de 5\,Bs y 5 billetes de 10\,Bs. ($x + y = 12$, $5x + 10y = 85$. Sustitución: $y = 12 - x$, $5x + 10(12 - x) = 85 \Rightarrow -5x = -35 \Rightarrow x = 7$.)

  \item $2{,}44$ horas ($\approx 2$\,h\,26\,min) después de que sale el segundo tren. (Cuando sale el 2do tren, el 1ro lleva $80 \times 2 = 160$\,km. Distancia restante: $440$\,km. Velocidad de acercamiento: $180$\,km/h. Tiempo: $440/180 = 22/9 \approx 2{,}44$\,h.)

  \item 17 correctas, 8 incorrectas. ($c + i = 25$, $4c - i = 60$. Sumando: $5c = 85 \Rightarrow c = 17$.)

  \item $\frac{12}{5} = 2{,}4$ horas (2\,h\,24\,min). (Tasa combinada: $\frac{1}{6} + \frac{1}{4} = \frac{5}{12}$. Tiempo: $\frac{12}{5}$\,h.)

  \item Luis: 22, Ana: 32, Pedro: 64 caramelos. (Sea $L$ la cantidad de Luis. Ana: $L + 10$. Pedro: $2(L + 10)$. Total: $L + (L + 10) + 2(L + 10) = 120 \Rightarrow 4L + 30 = 120 \Rightarrow L = 22{,}5$... Corrección: $L + (L+10) + 2(L+10) = 120$, $4L + 30 = 120$, $L = 22{,}5$. Si se requieren enteros, revisar enunciado. Respuesta exacta: Luis $= 22{,}5$, Ana $= 32{,}5$, Pedro $= 65$.)

  \item 8 días. Los primeros 7 días sube neto 1\,m/día (llega a 7\,m al inicio del día 8). El día 8 sube 3\,m y alcanza los 10\,m sin resbalar.

  \item 14 pruebas superadas. ($s + f = 20$, $5s - 3f = 52$. Sustituyendo $f = 20 - s$: $5s - 3(20 - s) = 52 \Rightarrow 8s = 112 \Rightarrow s = 14$.)

  \item Coloca una moneda en cada platillo de la balanza y deja la tercera fuera. Si la balanza se inclina, la moneda que sube es la falsa. Si queda equilibrada, la moneda fuera es la falsa.

  \item De 1 a 12: $1 + 2 + \cdots + 12 = \frac{12 \times 13}{2} = 78$ campanadas. En 24 horas se dan campanadas de 1 a 12 dos veces: $78 \times 2 = 156$. Además, se da 1 campanada cada media hora: hay $24 \times 2 = 48$ medias horas, pero 24 de ellas coinciden con las horas exactas (ya contadas). Quedan $48 - 24 = 24$ campanadas de media hora. Total: $156 + 24 = 180$ campanadas.
\end{enumerate}

\subsection*{Lenguaje Algebraico (Ejercicios 25--36)}

\begin{enumerate}[leftmargin=*, label=\textbf{\arabic*.}]
  \setcounter{enumi}{24}
  \item $2x + 9$ (donde $x$ es el número).

  \item $\frac{x}{2} = 20$ (donde $x$ es el número), es decir, $x = 40$.

  \item $n + (n + 1) + (n + 2) = 3n + 3$ (donde $n$ es el primero de los tres consecutivos).

  \item $x^2 - x = 12$ (donde $x$ es el número).

  \item Sea $h$ la edad del hijo. Pedro: $3h$. Ecuación: $h + 3h = 48$, es decir, $4h = 48$.

  \item $P = 3\ell$ (donde $\ell$ es la longitud del lado).

  \item Precio final $= P - 0{,}15P = 0{,}85P$.

  \item $v = \dfrac{d}{t}$ (donde $v$ = velocidad, $d$ = distancia, $t$ = tiempo).

  \item $(3 + 5) \cdot t = 40$, es decir, $8t = 40$ (donde $t$ es el tiempo en horas).

  \item $(a + b)^2 = a^2 + 2ab + b^2$.

  \item Ecuación 1: $d = v \cdot t$. Ecuación 2: $d = (v + 20)(t - 1)$.

  \item Sea $c$ el precio de una camisa. Pantalón: $2c$. Sistema: $3(2c) + 5c = 770$ y $2c = 2c + 40$... Corrección del segundo dato: «2 camisas cuestan 40\,Bs más que un pantalón» → $2c = 2c + 40$ no tiene sentido; la interpretación correcta es $2c - 2c = 40$. Re-lectura: $2c = p + 40$, con $p = 2c$, da $2c = 2c + 40$ (inconsistencia del enunciado). Sistema planteable: $6c + 5c = 770 \Rightarrow 11c = 770 \Rightarrow c = 70$. Pantalón: $140$\,Bs. Verificación con el segundo dato: $2(70) = 140$ y $140 + 40 = 180 \neq 140$. El sistema es sobredeterminado e inconsistente como ejercicio de planteamiento.
\end{enumerate}

\end{document}
